%%%%%%%%%%%%%%%%%%%%%%%%%%%%%%%%%%%%%%%%%
% "ModernCV" CV and Cover Letter
% LaTeX Template
% Version 1.11 (19/6/14)
%
% This template has been downloaded from:
% http://www.LaTeXTemplates.com
%
% Original author:
% Xavier Danaux (xdanaux@gmail.com)
%
% License:
% CC BY-NC-SA 3.0 (http://creativecommons.org/licenses/by-nc-sa/3.0/)
%
% Important note:
% This template requires the moderncv.cls and .sty files to be in the same
% directory as this .tex file. These files provide the resume style and themes
% used for structuring the document.
%
%%%%%%%%%%%%%%%%%%%%%%%%%%%%%%%%%%%%%%%%%

%----------------------------------------------------------------------------------------
%	PACKAGES AND OTHER DOCUMENT CONFIGURATIONS
%----------------------------------------------------------------------------------------

\documentclass[11pt,a4paper,sans]{moderncv} % Font sizes: 10, 11, or 12; paper sizes: a4paper, letterpaper, a5paper, legalpaper, executivepaper or landscape; font families: sans or roman

\moderncvstyle{classic} % CV theme
\moderncvcolor{blue} % CV color

\usepackage[scale=0.85, top=1.5cm, bottom=1.5cm]{geometry} % Optimized margins
\setlength{\hintscolumnwidth}{2.5cm} % Slightly narrower date column to give more room for text
\nopagenumbers{} % Ensure no page numbers interfere with space

%----------------------------------------------------------------------------------------
%	NAME AND CONTACT INFORMATION SECTION
%----------------------------------------------------------------------------------------
\usepackage{tikz}

\usepackage{mathptmx}
\usepackage[11pt]{moresize}

\firstname{San} % Your first name
\familyname{Cao} % Your last name
\title{Backend Engineer}

% All information in this block is optional, comment out any lines you don't need

\makeatletter
\@ifpackageloaded{moderncvstyleclassic}{%
\let\oldsection\section%
\renewcommand{\section}[1]{\leavevmode\unskip\vspace*{-\baselineskip}\oldsection{#1}}%
}{%
}
\makeatother

\email{caothesan@gmail.com}
\homepage{houtaru.github.io}
\social[linkedin]{thesancao}
\social[github]{houtaru}
%\homepage{https://github.com/priyankapps/} {https://github.com/priyankapps/}}

% \extrainfo{\faGithub\ \href{https://github.com/houtaru}{github.com/houtaru}}
%\href{https://github.com/houtaru}{ houtaru}}

\begin{document}

\makecvtitle
\vspace*{-1.5em}




%----------------------------------------------------------------------------------------
%	Experiments SECTION
%----------------------------------------------------------------------------------------

\section{Work Experience}

\cventry{May 2022 -- Present}{Software Engineer (Distributed Systems)}{Zalo}{}{}{
  \begin{itemize}
    \item Engineered \textbf{zero-downtime migration protocols} and automated cross-DC replication for \textbf{1,000+ in-house DB clusters}, enabling warm-standby DR.
    \item Designed metrics ingestion pipeline at \textbf{~1M samples/sec} with \textbf{~10s anomaly detection}; architected Thanos long-term store, cutting query latency 40\% and storage cost \textbf{3.5x}.
    \item \textbf{Encryption}: added in-memory encryption to Memcached, securing PII for 70M+ MAU with <2\% latency impact.
    \item \textbf{Backup \& Restore}: engineered daily \textbf{10TB+} backup flow with automated recoverability tests and snapshot drift metrics vs prod/previous.
    \item \textbf{Hermetic Build}: created reproducible build pipeline from source to binary, deployed to \textbf{2,000+ DB nodes} with zero-drift releases.
    \item \textbf{Stack}: C++, Java, NoSQL, Linux, Prometheus, Thanos, Grafana, Apache Thrift, Docker.
  \end{itemize}
}

\cventry{May 2021 -- Nov 2021}{Software Engineer Intern}{VNG Corporation}{}{}{
    \begin{itemize}
        \item Developed new features and maintained a card-games called Truco for Argentinians on server side.
        \item Maintained tool used for analysing user behavior based on given metric log.
        \item Technologies: Java, Memcached, Python, Hadoop, Spark, MySQL.
    \end{itemize}
}

\section{Skills}
\cvitem{}{
    \begin{itemize}
        \item \textbf{Languages:} C/C++, Python, Java, Go, SQL.
	\item \textbf{Technologies:} Linux, Docker, Prometheus, Thanos, Apache Thrift, Message Queue, NoSQL.
        \item \textbf{Core Technologies:} Distributed Systems, Data Structures/Algorithms, System Reliability Engineering.
    \end{itemize}
}

\section{Projects}
\cventry{2019 -- 2021}{Free Contest}{Problem Setter}{}{}{Collaborated with teams to engineer competitive programming problems for weekly contests on the self-hosted platform, focused on Vietnam National Olympiad in Informatics (VOI) preparation. Authored technical statements, robust testsets, and solution editorials.}
\cventry{2018}{cowboycoder.tech}{}{}{}{
A website sharing problem solution and code, as well as series of learning tips written or
translated in Vietnamese.}

%----------------------------------------------------------------------------------------
%	EDUCATION SECTION
%----------------------------------------------------------------------------------------

\section{Education}

\cventry{2018--2022}{Vietnam National University Ho Chi Minh City - University of Science}{}{}{}{\textit{Bachelor’s Degree in Computer Science, Advanced Program in Computer Science}\\ Published 2 research papers in AI (SRGAN \& Video Classification) at NAFOSTED NICS and MediaEval 2020.}

%\cventry{2015--2018}{Bao Loc High School for the Gifted}{}{}{}{}

%----------------------------------------------------------------------------------------
%	AWARDS SECTION
%----------------------------------------------------------------------------------------

\section{Honours and Awards}
\cventry{2021}{Top 10 in AIOT Innovation Works}{}{}{}{}
\cventry{2020}{Fourth prize in AI Challenge of AISocha 2020}{}{}{}{}
\cventry{2018}{Second Prize in ACM-ICPC Vietnam Southern Provincial Contest}{}{}{}{}
\cventry{2018}{Third Prize in Vietnamese National Olympiad in Informatics}{}{}{}{}

\end{document}



